% from main.tex
%%%%%%%%%%  scan idx 134  |  folio 121  %%%%%%%%%%
strictement parfait, la flèche \(\Id_F\) est localement cofinale dans \((\Qis/F)\).

\textbf{Notation 4.5.} Soit \(\mathcal M\) une \(S\)-catégorie fibrée. Soient \(X\in\Ob S\), et \(E,F\in\Ob\mathcal M_X\). Nous noterons \(\underline{\Hom}_{\mathcal M}(E,F)\) le faisceau sur \(S/X\) associé au préfaisceau des \(S\)-morphismes de \(E\) dans \(F\) dans \(\mathcal M\) (pour la définition de ce préfaisceau, voir [G] I 2.6; on prendra garde que, même si \(\mathcal C\) est un champ ([G]), le préfaisceau correspondant pour \(\mathcal M=\Kpx(\mathcal C)\) (ou \(\Dpx(\mathcal C)\)) n'est pas en général un faisceau).

\textbf{Corollaire 4.6.} Soient \(E,F\in\Ob\Cpx(\mathcal C_S)\), \(E\) étant strictement parfait. La flèche canonique
\[
  \underline{\Hom}_{\Kpx(\mathcal C)}(E,F)
  \longrightarrow
  \underline{\Hom}_{\Dpx(\mathcal C)}(E,F)
\]
est un isomorphisme.

\emph{Preuve.} C'est une conséquence immédiate de (4.4) et de la formule
\[
  \underline{\Hom}_{\Dpx(\mathcal C_X)}(E_X,F_X)
  = \varinjlim_{(\Qis/E_X)^\circ}\underline{\Hom}_{\Kpx(\mathcal C_X)}(\,\cdot\,,F_X)
\]
pour \(X\) au-dessus de \(S\).

\textbf{Définition 4.7.} Soit \(F\in\Ob\Dpx(\mathcal C_S)\) et soient \(a,b\in\mathbb Z\) tels que \(a\leq b\). Nous dirons que \(F\) est d'amplitude \(\mathcal C_0\)-parfaite contenue dans l'intervalle \([a,b]\), et nous écrirons
\[
  \mathcal C_0\text{-}\Parf\text{-}\amp(F)\subset[a,b],
\]
si l'ensemble des \(X\) au-dessus de \(S\) tels qu'il existe un quasi-isomorphisme \(F'\to F\) de \(\Cpx(\mathcal C_X)\)
